% ===================================================================
%  LENTELĖ Nr. 13 — Viešbutis. --- Restoranas. Kavinė.
%  Traduction 2026 d'apres texto/io, controlee sur texto/fr, texto/ru
%  et texto/cs. Consigne : voir texto/lt/10-lentele-01.tex.
%
%  LE POURBOIRE DE L'ALINEA 9 EST LE MOT LE PLUS INSTRUCTIF DE LA
%  COLONNE, ET IL FAUT TROIS LANGUES POUR LE LIRE. Le tcheque dit
%  « spropitné », calque de l'allemand « Trinkgeld » : l'argent pour
%  BOIRE. Le lituanien dit « arbatpinigiai », l'argent pour LE THE,
%  qui n'est pas le calque de l'allemand mais celui du russe
%  « чаевые ». Le francais et l'ido disent le pourboire et la
%  gratifiketo.
%
%  TROIS LANGUES CALQUENT LE MEME MOT, ET CHACUNE CALQUE SON PROPRE
%  VOISIN. Le calque, troisieme voie enseignee par la colonne
%  tcheque, ne dit donc pas seulement qu'une langue a emprunte une
%  idee : IL DIT DE QUI. « Boire » designe Berlin, « the » designe
%  Moscou, et l'on n'a pas besoin d'histoire pour le savoir — le mot
%  le porte.
%
%  ET C'EST LA MEME LIGNE QUE CELLE DU TABLEAU 11, PRISE PAR
%  L'AUTRE BOUT. La ville lituanienne avait refait la caisse
%  d'epargne et les pompiers contre le polonais et le russe ; le
%  cafe lituanien garde du russe sa maniere de nommer le pourboire.
%  Une langue qui a decide contre un voisin ne s'est pas defaite de
%  lui.
%
%  LE RESTE DU TABLEAU SUIT LES REGLES DEJA VUES. « viešbutis »
%  l'hotel est un compose — « viešas » public et « butas » logement ;
%  « kavinė » le cafe et « valgykla » le restaurant portent le
%  suffixe de lieu du tableau 6 ; « padavėjas » le garcon est celui
%  qui donne et « virėjas » le cuisinier celui qui fait bouillir.
%  Restent etrangers « restoranas », « biliardas », « omnibusas »,
%  « telegrama », « domino », « triktrakas » : les jeux et les
%  machines, comme au tableau 12.
%
%  QUATRE GENITIFS TENUS PAR UNE RELATIVE OU UN PARTICIPE : le
%  cocher (30) de l'omnibus (29), la malle (31) du voyageur (32), les
%  nuages (67) de l'horizon (66), l'enseigne (69) du hall (68).
% ===================================================================

%%K t13-noto-1 noto
\VUnotes{143.2mm}{%
\textit{\VUgras{(*) Dėl kambario smulkmenų žiūrėk lentelę Nr.~6.}}%
}

%%K t13-apar-1 apar
\VUsaut{3.0mm}
\VUtitre{15.0pt}{60}{LENTELĖ Nr. 13}
\VUsaut{1.6mm}
\VUpk{10.2pt}{40}{Viešbutis. --- Restoranas.}
\VUsaut{2.0mm}
\VUpk{10.2pt}{40}{Kavinė.}
\VUsaut{3.4mm}

%%K t13-01-1 p
1. --- Vakar vakare atvykęs į Ruajaną, apsistojau „\textit{Vandenyno viešbutyje}“,
kurį man buvo rekomendavęs draugas.

%%K t13-01-2 p
Man davė gražų \VUgras{kambarį} \textsuperscript{(2)} trečiajame \VUgras{aukšte}
\textsuperscript{(1)}, kur labai gerai miegojau (*).

%%K t13-02-1 p
2. --- Šįryt, išeidamas iš savo kambario, į kurį \VUgras{tarnaitė}
\textsuperscript{(3)} buvo atnešusi pusrytėlius, įėjau į \VUgras{rūkomąjį}
\textsuperscript{(5)}, kuris naudojamas ir kaip \VUgras{skaitykla}
\textsuperscript{(4)}, kad paskaityčiau \VUgras{laikraščius} \textsuperscript{(6)},
rūkydamas \VUgras{cigaretę} \textsuperscript{(8)}. Paskaitęs nusileidau
\VUgras{liftu} \textsuperscript{(9)} ir nuėjau pasivaikščioti po miestą.

%%K t13-03-1 p
3. --- Kadangi buvau pamiršęs paklausti viešbučio \VUgras{tarnautojo}
\textsuperscript{(12)}, kurią valandą duodami valgiai prie \VUgras{bendro stalo}
\textsuperscript{(11)}, grįžau anksti ir pasinaudojau proga išsimaudyti viešbučio
\VUgras{vonioje} \textsuperscript{(10)}. Paskui vėl nuėjau prie \VUgras{kasos}
\textsuperscript{(15)}, kad paskambinčiau vienam savo draugui. Bet \VUgras{telefonas}
\textsuperscript{(13)} buvo užimtas; matydama mano sutrikimą, \VUgras{kasininkė}
\textsuperscript{(14)}, kuri įrašinėjo \VUgras{sąskaitą} \textsuperscript{(17)} į
\VUgras{svečių knygą} \textsuperscript{(16)}, paprašė \VUgras{durininko}
\textsuperscript{(18)}, kad pasiųstų \VUgras{pasiuntinuką} \textsuperscript{(19)}
atlikti mano pavedimo \VUgras{dviračiu} \textsuperscript{(20)}.

%%K t13-04-1 p
4. --- Padėkojau jai ir išėjau duoti arbatpinigių \VUgras{nešikui}
\textsuperscript{(24)} (juodadarbiui), kuris savo \VUgras{rankiniu vežimėliu}
\textsuperscript{(25)} atvežė iš stoties mano \VUgras{skrybėlių dėžę}
\textsuperscript{(26)}, \VUgras{mano lagaminą} \textsuperscript{(27)} ir \VUgras{mano
kelionmaišį} \textsuperscript{(28)}. Ką tik atvykęs \VUgras{laiškininkas}
\textsuperscript{(21)} padavė man \VUgras{laišką} \textsuperscript{(22)}.

%%K t13-05-1 p
5. --- Dar kurį laiką pastovėjau ant durų slenksčio, prie \VUgras{pašto dėžutės}
\textsuperscript{(23)}, žiūrėdamas į judėjimą gatvėje. \VUgras{Vežikas}
\textsuperscript{(30)}, kuris valdo \VUgras{omnibusą} \textsuperscript{(29)}, krovė ant
savo vežimo \VUgras{lagaminą} \textsuperscript{(31)}, priklausantį \VUgras{keleiviui}
\textsuperscript{(32)}, kurį lydėjo \VUgras{viešbučio šeimininkas}
\textsuperscript{(33)}. Jaunas \VUgras{telegrafistas} \textsuperscript{(35)} neskubėdamas
atnešė \VUgras{telegramą} \textsuperscript{(34)} viešbučio klientui;
\VUgras{turistas} \textsuperscript{(36)} su savo \VUgras{fotoaparatu}
\textsuperscript{(38)} per petį žiūrėjo į savo \VUgras{kelionvedį}
\textsuperscript{(37)}.

%%K t13-tit-1 sub
\VUsaut{4.0mm}
\VUpk{10.2pt}{40}{Pusryčių metas.}
\VUsaut{3.0mm}

%%K t13-06-1 p
6. --- Prieš grotas nerūpestingas \VUgras{pasiuntinys} \textsuperscript{(39)}
snūduriavo tarp savo \VUgras{nešimo kablio} \textsuperscript{(40)} ir savo \VUgras{batų
tepimo dėžės} \textsuperscript{(41)}, o jauna \VUgras{skalbėja} \textsuperscript{(42)},
kuri nešė skalbinių \VUgras{krepšį} \textsuperscript{(43)}, juokėsi iš iškilmingos
\VUgras{metrdotelio} \textsuperscript{(44)} minos, kuris stovėjo prie durų.

%%K t13-07-1 p
7. --- Matydamas \VUgras{virėją} \textsuperscript{(45)}, kuris išėjo iš savo
\VUgras{virtuvės} \textsuperscript{(47)}, esančios \VUgras{rūsio aukšte}
\textsuperscript{(48)}, kad pasiųstų \VUgras{indų plovėją} \textsuperscript{(46)}
pavedimo, prisiminiau, kad pusryčių valanda arti, ir įėjau į valgyklą, eidamas per
kiemą, kuriame \VUgras{vairuotojas} \textsuperscript{(50)} statė savo
\VUgras{automobilį} \textsuperscript{(49)}, o \VUgras{mechanikas}
\textsuperscript{(52)} taisė, pagal \VUgras{dviratininko} \textsuperscript{(53)}
nurodymus, \VUgras{benzininį triratį} \textsuperscript{(51)}, kuris ką tik buvo
sudaužytas.

%%K t13-08-1 p
8. --- Paklausiau jauno \VUgras{pasiuntinuko} \textsuperscript{(54)}, kuris nešė
\VUgras{alaus stiklines} \textsuperscript{(55)}, ar bendras stalas yra po veranda, ir,
jam patvirtinus, užėmiau vietą prie vieno iš stalų ir liepiau paduoti man puikų
valgį.

%%K t13-noto-2 noto
\VUnotes{143.2mm}{%
\textit{\VUgras{(*) Pasitelkdamas į žodyną įtrauktą patiekalų sąrašą, mokytojas
galės lengvai įvairinti žemiau pateiktą valgiaraštį.}}%
}

%%K t13-tit-2 sub
\VUsaut{4.0mm}
\VUpk{10.2pt}{40}{Puikūs pusryčiai.}
\VUsaut{3.0mm}

%%K t13-09-1 p
9. --- Padavėjas atnešė man pusryčių valgiaraštį (*) ir vynų sąrašą. Išsirinkau ir
užsisakiau kaip \VUgras{užkandį krevečių} su \VUgras{sviestu,} o kaip pirmą
patiekalą --- \VUgras{Kano būdu troškintus vidučius.} \VUgras{Žuvies} nevalgiau,
bet paprašiau porcijos ėrienos \VUgras{kumpio} ir kaip \VUgras{daržovių patiekalo}
--- \VUgras{špinatų} su grietinėle. Paskui, kadangi nė vienas iš \VUgras{tarpinių patiekalų}
man nepatiko, liepiau atnešti įvairių desertų: \VUgras{sūrio,} \VUgras{vaisių} ir
\VUgras{sausainių.} Dar pridėjau puikaus Saint-Emilion \VUgras{vyno} ir, labai
patenkintas savo valgiu, palikau stalą, padovanojęs \VUgras{padavėjui}
\textsuperscript{(57)} gražių \VUgras{arbatpinigių} \textsuperscript{(59)}.

%%K t13-10-1 p
10. --- Matydamas mane pasirengusį išeiti, \VUgras{valdytojas} \textsuperscript{(60)}
pasiūlė man užeiti į balkoną pasigėrėti puikia \VUgras{Vandenyno}
\textsuperscript{(62)} panorama. Toli buvo matyti žvejų \VUgras{valtelės}
\textsuperscript{(63)}, \VUgras{burlaiviai} \textsuperscript{(64)} ir milžiniškas
\VUgras{keleivinis garlaivis} \textsuperscript{(65)}, kurio juodas siluetas išsiskiria
baltuose \VUgras{debesyse} \textsuperscript{(67)} ties \VUgras{horizontu}
\textsuperscript{(66)}. Lengvas vėjelis plaikstė \VUgras{vėliavėlę}
\textsuperscript{(70)}, įsmeigtą virš \VUgras{iškabos} \textsuperscript{(69)}, kuri
kabo ant \VUgras{halės} \textsuperscript{(68)}.

%%K t13-tit-3 sub
\VUsaut{3.0mm}
\VUpk{10.2pt}{40}{Pramogos.}
\VUsaut{3.0mm}

%%K t13-11-1 p
11. --- Reginys buvo toks įdomus, kad nusprendžiau rytoj užlipti į kavinės terasą,
pasiimdamas \VUgras{žiūronėlį} \textsuperscript{(71)} ar \VUgras{žiūroną}
\textsuperscript{(72)}.

%%K t13-12-1 p
12. --- Palikęs balkoną, nuėjau į viešbučio kavinę išgerti \VUgras{gėrimo}
\textsuperscript{(73)} ir sužaisti \VUgras{biliardo partiją} \textsuperscript{(75)}.
Nesustodamas perėjau kavos salę, kurioje daug lankytojų žaidė \VUgras{šachmatais}
\textsuperscript{(78)}, \VUgras{šaškėmis} \textsuperscript{(79)}, \VUgras{domino
kauliukais} \textsuperscript{(80)}, \VUgras{triktraku} \textsuperscript{(81)} ar
\VUgras{kortomis} \textsuperscript{(82)}, ir tiesiai užlipau į \VUgras{biliardinę}
\textsuperscript{(74)}.

%%K t13-13-1 p
13. --- Kai partija baigėsi, nusileidau į pirmąjį aukštą ir paprašiau padavėjo
\VUgras{voko} \textsuperscript{(84)}, \VUgras{popieriaus} \textsuperscript{(83)},
\VUgras{pašto ženklo} \textsuperscript{(85)}, \VUgras{plunksnos}
\textsuperscript{(87)} ir \VUgras{rašalo} \textsuperscript{(86)} laiškui parašyti.

%%K t13-13-2 p
Paskui pasišaukiau \VUgras{vežiką} \textsuperscript{(92)}, kurio \VUgras{ekipažas}
\textsuperscript{(88)} buvo sustojęs prieš kavinę. Paklausiau jo kainos pagal valandas
ar už vieną kursą, ir, kai susitarėme dėl kainos, jis atidarė man \VUgras{ekipažo
dureles} \textsuperscript{(89)} ir įsodino mane. Jis vėl užlipo ant \VUgras{pasostės}
\textsuperscript{(91)}, greitai pajudino arklius \VUgras{botagėliu}
\textsuperscript{(93)}, ir mes išvykome žavingai pasivažinėti.

\VUsaut{6.0mm}
\VUfilet{22mm}
